\section{Nonconvex quadratic and algebraic cone arithmetic}
\label{app:further-arithmetic}

Appendix~\ref{app:qc} distinguishes native convex quadratics from squared
cone inequalities. This appendix treats two further arithmetic interfaces.
For arbitrary weak quadratic inequalities, a fixed Hessian span gives short
algebraic feasible points and finite-infimum descriptions, even when the
infimum is not attained. Finding these descriptions remains hard. For actual
second-order cone systems, convexity permits a polynomial-time decision and
witness algorithm at fixed integer dimension and fixed continuous Hessian
span, including coefficients in one explicitly represented number field.
The final linear optimization calls have rational coefficients.

\subsection{Weak quadratic systems and a sampling corollary}
\label{app:ncq-setting}

Let
\begin{equation}
\label{eq:ncq-system}
 S=\{x\in\R^n:Ax\le b,\ Ex=e,\ g_i(x)\le0\ (1\le i\le m)\},
 \qquad g_i(x)=\tfrac12x^{\mathsf T}Q_ix+a_i^{\mathsf T}x+c_i,
\end{equation}
where the symmetric matrices $Q_i$ may be indefinite. Unless a field is
specified, all coefficients are rational and explicitly encoded. Set
$h=\dim_{\Q}\Span\{Q_1,\ldots,Q_m\}$. An arbitrary quadratic objective
$f$ is included in the total binary input length $L$, but its Hessian is
excluded from $h$. All inequalities are weak. There is no assumed box,
smoothness, convexity, or constraint qualification.

We also use a coefficient-sensitive accounting convention. Let $S_0\ge2$
bound the numbers of variables, rows, coefficient positions, and index bits,
and let $\tau_0\ge1$ bound each rational numerator and denominator bit
length. The bounds below have the form
\begin{equation}
\label{eq:ncq-budget}
 \text{degree}\le S_0^{C(h+1)},\qquad
 \text{coefficient bits and total representation length}
       \le(\tau_0+1)S_0^{C(h+1)},
\end{equation}
for an effective absolute constant $C$, enlarged when necessary. In
particular they are $L^{O(h+1)}$. The separate accounting prevents a large
printed radius from being counted as a large structural dimension.

Choose a Hessian basis $B_1,\ldots,B_h$ among the $Q_i$ and solve
$Q_i=\sum_j\gamma_{ij}B_j$. Introduce
\begin{equation}
\label{eq:ncq-lift}
 y_j=\tfrac12x^{\mathsf T}B_jx,\quad w=(x,y),\quad
 F_j(w)=\tfrac12x^{\mathsf T}B_jx-y_j.
\end{equation}
Then $S$ is in bijection with
$P\cap Z(F_1,\ldots,F_h)$, where $P$ retains every original affine row
and replaces each $g_i\le0$ by
$\sum_j\gamma_{ij}y_j+a_i^{\mathsf T}x+c_i\le0$.
Rational elimination and determinant bounds give polynomial structural
size and coefficient bits $(\tau_0+1)S_0^{O(1)}$. Every consistent
subset of the rows of $P$ taken as equalities has a chart
\begin{equation}
\label{eq:ncq-chart}
 w=a+Vu,\qquad u\in\R^d,
\end{equation}
where $V$ has full column rank and the same coefficient-bit bound. Choose
independent normals and free coordinates to obtain this chart by Cramer's
rule. There are finitely many charts, at most $2^{S_0^{O(1)}}$; no efficient
enumeration is asserted.

\begin{lemma}[A component in a polyhedron]
\label{lem:ncq-face}
If a closed polyhedron $P$ meets a real algebraic set $V_0$, some face $H$
of $P$ has a connected component of $V_0\cap\operatorname{aff}H$ wholly
contained in $H$.
\end{lemma}

\begin{proof}
Choose a face $H$ of minimum dimension among those meeting $V_0$, allowing
$P$ itself. Every relative-boundary point of $H$ belongs to a proper face
of $H$, also a face of $P$. Thus
$V_0\cap\operatorname{relbd}H=\varnothing$. Choose a connected component
$C$ of $V_0\cap\operatorname{aff}H$ meeting $H$. The nonempty set
$C\cap H$ is closed in $C$ and equals $C\cap\operatorname{relint}H$,
which is open in $C$. Hence $C\subseteq H$. The argument includes
unbounded faces, lineality, and affine-space polyhedra.
\end{proof}

\begin{theorem}[Algebraic feasible points]
\label{thm:ncq-feasible}
If the rational system~\eqref{eq:ncq-system} is nonempty, it contains a
point in one number field with degree and total common-field representation
length $L^{O(h+1)}$. The coefficient-sensitive bounds
\eqref{eq:ncq-budget} hold. Consequently feasibility is in $\NP$ for fixed
$h$, and a feasible point exists in a box of radius $2^{L^{O(h+1)}}$.
\end{theorem}

\begin{proof}
Apply Lemma~\ref{lem:ncq-face} to the lift~\eqref{eq:ncq-lift}. A chart
of the selected face has bounded-size rational coefficients. A
zero-dimensional chart directly gives a rational feasible point. Otherwise,
in that chart,
Grigoriev--Pasechnik's component-sampling theorem
\cite[Theorem~1.2]{GrigorievPasechnik2005}, applied to the quadratic map
$(F_1,\ldots,F_h)$ and the outer polynomial $\sum_jY_j^2$, gives a
univariate sample in every real connected component. Their theorem allows
degenerate and unbounded sets; for integer input, its degrees and output
coefficient bits are bounded by $(2d)^{O(h)}$ times the input coefficient
size, with a harmless polynomial factor for all coordinates. Clear the
restricted rational denominators first. One sample lies in the component
of Lemma~\ref{lem:ncq-face}, so it satisfies every row of $P$.

The samples use one real algebraic root and rational coordinate functions
with denominator nonzero there. Select its irreducible factor, invert the
coordinate denominator modulo that factor, and isolate the selected real
root. Dense univariate arithmetic preserves~\eqref{eq:ncq-budget}; the
integer-factor bound~\eqref{eq:qc-minpoly-height} controls the factor's
coefficients. The affine chart gives the required common-field point. For
$h=0$ the conclusion is the ordinary bounded-size rational feasible-point
property of polyhedra, also obtained by minimizing the squared norm and
solving an independent active affine chart.

The verifier receives only this represented point. It checks its root and
all original row signs by Lemma~\ref{lem:qc-sign}; it does not verify a
choice of face or a sampling transcript. This proves the $\NP$ assertion.
Cauchy's bound on the coordinate polynomials gives the stated radius.
\end{proof}

This theorem is an explicit polyhedral corollary of established sampling.
It does not import the optimization statement announced as Theorem~1.5 in
\cite{GrigorievPasechnik2005}, whose proof was deferred there. Optimization
requires the separate limit argument below.

Deleting inactive rows would give an incorrect global argument. On the
connected ellipse $4(x-2)^2+y^2=4$, the squared norm equals
$-3x^2+16x-12$. With $x\ge5/2$, its minimum is $9$ at $(3,0)$, while
the other endpoint has value $37/4$. The affine row is inactive at $(3,0)$.
Deleting it exposes $(1,0)$, of squared norm $1$. We retain all original
rows throughout and use active charts only for local stationarity.

\subsection{Generic perturbations and bounded algebraic outputs}
\label{app:ncq-elimination}

The next argument also works over a selected real embedding of a number
field $K$. It will be used for squared cone rows, whose Hessians can be
indefinite. Explicit field input means a dense primitive irreducible
polynomial, a selected real-root isolator, and rational power-basis
coefficient lists, as in~\eqref{eq:cone-field} below. Its degree is
counted in $L$.

\begin{lemma}[Generic quadratic critical equations]
\label{lem:ncq-generic}
Let $r,f_1,\ldots,f_s$ be quadratics in $d>0$ complex variables, with
all their coefficients allowed to vary. If $s\le d$, outside a proper
algebraic set every common zero of the $f_i$ has independent gradients,
and at every KKT root both
\[
 M=\nabla^2r+\sum_i\lambda_i\nabla^2f_i,
 \qquad \begin{pmatrix}M&G^{\mathsf T}\\G&0\end{pmatrix}
\]
are nonsingular, where $G$ has rows $\nabla f_i^{\mathsf T}$. If $s>d$,
generic $f_i$ have no common zero. The exceptional set is contained in the
zero set of a nonzero polynomial of degree $2^{\operatorname{poly}(d+s)}$.
No coefficient-height bound for that polynomial is required.
\end{lemma}

\begin{proof}
At fixed $u$, constants and linear terms prescribe each $f_i(u)$ and
$\nabla f_i(u)$ independently. Vanishing values have codimension $s$;
rank deficiency of an $s\times d$ matrix has codimension $d-s+1$.
Including the $d$ coordinates of $u$ leaves bad-incidence dimension at most
one less than coefficient-space dimension. For $s>d$, the values alone
give this strict inequality.

For KKT roots, solve the constraint equations for their constant terms and
stationarity for the objective's linear coefficients. The incidence is an
affine space of the same dimension as coefficient space. Neither displayed
determinant vanishes identically on it. Indeed, take
\[
 f_i=u_i^2-1\ (i\le s),\qquad
 r=\sum_{j=1}^du_j^2+\sum_{i=1}^su_i.
\]
At $u_i=\pm1$ for $i\le s$, $u_j=0$ for $j>s$, and
$\lambda_i=-1-1/(2u_i)$, the diagonal entries of $M$ are $-1/u_i$ and
$2$, respectively. The rows of $G$ are $2u_ie_i^{\mathsf T}$, so
$GM^{-1}G^{\mathsf T}$ is diagonal and invertible. The determinant-zero
incidences therefore have dimension at most one less than coefficient
space; their projected closures are proper.

For the quantitative bound, describe gradient dependence by a nonzero
dependence vector with one coordinate normalized to one. The $s$ such
charts use polynomially many equations of degree at most three. The other
incidences use the value and stationarity equations and one determinant,
again polynomially many variables and polynomial degrees. Use cumulative
affine degree, summing the degrees of all irreducible components. Affine
B\'ezout bounds this degree by $2^{\operatorname{poly}(d+s)}$, and linear
projection does not increase it
\cite[Section~1.2.1]{KrickPardoSombra2001}. A proper irreducible image
component of dimension $r_0$ has a generic linear projection into $r_0+1$
coordinates whose image closure is a hypersurface of no greater degree;
pulling back its equation gives a containing hypersurface. Multiplying
these equations over all components gives the asserted bound. These bounds
use cumulative degree, so all components are charged.
\end{proof}

Both determinants are needed. For example, $M=\diag(1,-1)$ and $G=(1,1)$
are individually of full rank but $GM^{-1}G^{\mathsf T}=0$.

Consider the family on the original lifted polyhedron $P$,
\begin{equation}
\label{eq:ncq-perturbation}
 r_{\varepsilon,\eta}(w)=f(x)+\varepsilon\norm x^2+\eta P_0(w),
 \qquad |F_j(w)+\eta^2P_j(w)|\le\eta,
\end{equation}
where $P_0,\ldots,P_h$ are ambient quadratics. The one-parameter version
omits $\varepsilon\norm x^2$ and uses an arbitrary fixed quadratic
objective. For small-point proofs it uses $\norm w^2$ as that objective.

\begin{lemma}[One small perturbation tuple]
\label{lem:ncq-perturbation}
There are integer $P_j$ with coefficient bits $S_0^{O(1)}$ such that,
on every original affine chart and every subset of bands with one
orientation per index, the conclusions of Lemma~\ref{lem:ncq-generic}
hold for all sufficiently small positive $\eta$ at each fixed positive
$\varepsilon$ outside a finite exceptional set. The tuple is independent
of both parameters and of any auxiliary box. In the one-parameter version
the conclusions hold on one sufficiently small positive tail.
\end{lemma}

\begin{proof}
Restriction of ambient quadratics to a full-rank affine chart is onto the
space of chart quadratics: extend its affine coordinate functions to the
ambient space. At any fixed $\varepsilon$ and nonzero $\eta$, the factors
$\eta$ and $\eta^2$ in~\eqref{eq:ncq-perturbation} therefore give an
onto coefficient map for the objective and each selected oriented band.
Pulling back a nonzero exceptional polynomial produces a nonzero polynomial
in the perturbation coefficients and $(\varepsilon,\eta)$. Choose one
nonzero coefficient in its formal parameter expansion. The product over
all charts and oriented subsets is nonzero and has degree at most
$2^{S_0^{O(1)}}$.

A polynomial of total degree $D_0$ cannot vanish on all of
$\{0,\ldots,D_0\}^p$: induction on $p$ reduces this to the univariate
root bound. Thus the product has a nonzero value at an integer tuple with
$S_0^{O(1)}$ bits per coefficient. This reasoning holds over $K$ as well
as over $\Q$. Once the tuple is fixed, the exceptional polynomials in
$(\varepsilon,\eta)$ are nonzero. Only finitely many $\varepsilon$ make
any one identically zero in $\eta$, since they are roots of a nonzero
coefficient polynomial. For every other fixed $\varepsilon$, exclude its
finitely many positive $\eta$ roots. In one parameter only the last step
is needed. The permitted tail may depend on $\varepsilon$.
\end{proof}

At a perturbed minimizer where auxiliary box rows are inactive, make all
active original affine rows equalities and use~\eqref{eq:ncq-chart}. The
other affine rows are locally strict. Only one orientation per band can be
active, since the band has width $2\eta>0$. There are therefore
$s\le\min(h,d)$ active nonlinear rows. Lemma~\ref{lem:ncq-generic}
gives ordinary KKT necessity with objective multiplier one. Write the
restricted stationarity equations as $Mu+\mu=0$ and set
\[
 \Delta=\det M,\qquad \rho=-\operatorname{adj}(M)\mu,
 \qquad u=\rho/\Delta.
\]
For each active quadratic $\phi_i$, define
$H_i=\Delta^2\phi_i(\rho/\Delta)$. At its selected root,
\begin{equation}
\label{eq:ncq-jacobian}
 \frac{\partial H}{\partial\lambda}
       =-\Delta^2GM^{-1}G^{\mathsf T}.
\end{equation}
Indeed, differentiated stationarity gives
$\partial u/\partial\lambda_i=-M^{-1}\nabla\phi_i(u)$; terms
differentiating $\Delta^2$ vanish since $\phi_i(u)=0$.
The bordered determinant makes~\eqref{eq:ncq-jacobian} nonsingular by
the Schur complement. Neither positive definiteness nor bounded multipliers
is used.

Coordinates, fixed linear forms, squared norms, and quadratic objective
values have rational output formulas with common denominator $\Delta^2$.
The reduced equations and output numerators have multiplier degree
$a_0=S_0^{O(1)}$ and polynomial parameter degree. For rational input,
their full coefficient-norm logarithm is
$(\tau_0+1)S_0^{O(1)}$. To see the coefficient accounting, clear the
polynomially many input denominators once, before determinant expansion.
All later denominators divide bounded powers of this one denominator.
Determinant and adjugate coefficient sums multiply by at most factorial
and polynomial-size factors. Their logarithms have the stated bound even
when the expanded monomial list is large. The parameters remain formal;
no numerical parameter precision or unknown box length enters this bound.

Over an explicitly represented field $K$ of degree $D$, use the local
coefficient sums at archimedean places and maxima at finite places from
Lemma~\ref{lem:qc-elimination}. Field linear algebra has polynomial bit
cost: multiplication by a field element is a rational $D$-dimensional
linear map, and an invertible $r\times r$ field system is a rational
$rD$-dimensional system. The root height of the input generator is
bounded by its polynomial's coefficient height plus $\log2$. An encoded
coefficient $\sum_{j<D}a_j\alpha^j$ has height at most
$\sum_j\height(a_j)+(D-1)\height(\alpha)+\log D$.
Lemma~\ref{lem:qc-heights} therefore bounds input and chart scalar
heights polynomially. If $U_v\ge1$ bounds the restricted scalar
coefficients, each polynomial entry in $M$ and $\mu$ has local norm at
most $c(s+2)U_v$ at archimedean places, for an absolute $c$, and at
most $U_v$ at finite places. Its determinants and adjugate numerators
thus have norms bounded by factorials times polynomially many powers
of these bounds. A fixed number of these substitutions bounds the
averaged logarithmic local norm by
$E=L^{O(1)}$; it does not sum the heights of expanded monomials.

Consequently Lemma~\ref{lem:qc-elimination}, with its inner parameter
renamed $\eta$, applies to finite ordered output limits. Its bounds are
\begin{equation}
\label{eq:ncq-elimination-budget}
 \Lambda=(a_0+1)^s,\qquad
 W=\Lambda\bigl(a_0(s+1)+1\bigr)E+\log(\Lambda!)+\Lambda\log3.
\end{equation}
It supplies a nonzero polynomial over $K$ of degree at most $\Lambda$
and affine coefficient height at most $W$. A one-parameter application
uses the lemma with a dummy inner parameter and the same root at every
value of that parameter. Its proof requires nonsingular multiplier roots,
a nonzero output denominator, and finite ordered scalar limits; it has no
convexity or bounded-multiplier hypothesis. In its proof the lowest inner
parameter coefficient is extracted first, and the lowest outer parameter
coefficient second. Specializations making an intermediate relation
identically zero are permitted. These facts will matter for nonattainment.

\begin{lemma}[From one tuple to one short field]
\label{lem:ncq-common-field}
Suppose one selected limiting tuple $w^*$ has the preceding elimination
relations for all coordinates and every fixed $K$-linear form, with the
same degree bound $\Lambda$. Then $[K(w^*):K]\le\Lambda$. For explicit
field input the absolute field $\Q(\alpha,w^*)$ has degree and total
common-field representation length $L^{O(h+1)}$. For rational input the
coefficient-sensitive bounds~\eqref{eq:ncq-budget} hold.
\end{lemma}

\begin{proof}
Coordinate relations first make $K(w^*)/K$ finite and separable. A
primitive element is a $K$-linear combination of its generating
coordinates. Applying the same output relation to that combination proves
the relative degree bound. Different coordinates must use the same tuple;
multiplying their separate degrees would not give this conclusion.

A root $\xi$ of a degree-$e$ polynomial over $K$ with affine coefficient
height $W$ has absolute degree at most $De$ and height at most $W+\log2$.
The latter follows by the placewise root bounds and the product formula
for the leading coefficient, exactly as in the proof of
Theorem~\ref{thm:qc-height}. If $p_\xi$ is its primitive integer minimal
polynomial, the identity between Mahler measure and absolute
height~\cite{BombieriGubler2006}, followed by expansion in the roots, gives
\[
 \log\norm{p_\xi}_\infty
       \le De\bigl(W+2\log2\bigr).
\]
Thus every coordinate and the input generator $\alpha$ have absolute
minimal-polynomial coefficient bits $H=L^{O(h+1)}$.

Let $J=D\Lambda$ and apply the embedding argument to the generating tuple
$(\alpha,w^*)$. Each pair of distinct embeddings excludes one proper
hyperplane of integer-combination coefficients. Their product has degree
at most $J(J-1)/2$, so an integer combination with coefficients between
$0$ and $J(J-1)/2$ is primitive. Let $r_0$ be the tuple length and
$\mathfrak A$ the product of its coordinate minimal polynomials'
absolute leading coefficients. Then $\log_2\mathfrak A\le r_0H$,
and $\mathfrak A$ times that combination is integral. Every conjugate
of the combination is bounded by $r_0J^2 2^{H+1}$. Expanding its
degree-at-most-$J$ conjugate product after this integral scaling bounds
minimal-polynomial coefficient bits by
$O(J(r_0H+H+\log(r_0+1)+\log(J+1)))$.

For completeness, power-basis coordinates have equally controlled bits.
Scale the primitive element and one coordinate by their respective leading
minimal-polynomial coefficients, obtaining algebraic integers $A,B$.
Their conjugates have magnitude at most $2^{O(H'+1)}$ for the established
coefficient bound $H'$. If $d_0$ is the joint degree, write
$B=\sum_{r<d_0}c_rA^r$ and take traces after multiplication by
$A^s$, $0\le s<d_0$. The matrix
$(\operatorname{Tr}(A^{r+s}))_{r,s}$ and right side
$(\operatorname{Tr}(BA^s))_s$ are integers with
$O(d_0(H'+1)+\log d_0)$ bits. Separability makes the trace pairing
nonsingular. Cramer's rule bounds each rational $c_r$ by
$O(d_0^2(H'+1)+d_0\log d_0)$ bits. Scaling back gives all rational
coordinate maps. Discriminant root separation gives an isolator of
polynomial length in degree and coefficient bits. All these operations
multiply the previous bounds by polynomials in the degree and structural
size. For $K=\Q$ they preserve linear dependence on $\tau_0+1$, proving
\eqref{eq:ncq-budget}.
\end{proof}

\begin{lemma}[Nonconvex radius and boxed-value bounds over a field]
\label{lem:ncq-field-bounds}
For an explicitly represented selected real field $K$, a nonempty weak
quadratic system of Hessian span $h$ has a common-field feasible point of
degree and total length $L^{O(h+1)}$, hence one of magnitude at most
$2^{L^{O(h+1)}}$. On a supplied rational box, a nonempty such system has
an optimal point for any quadratic objective with the same bounds. Its
value has an integer annihilator of degree and coefficient bits
$L^{O(h+1)}$. A nonzero value has absolute value at least
$2^{-L^{O(h+1)}}$.
\end{lemma}

\begin{proof}
Use the lift over $K$ and the one-parameter perturbations of
Lemma~\ref{lem:ncq-perturbation}. For a feasible point, minimize
$\norm w^2+\eta P_0(w)$ on the unchanged polyhedron and the outward
bands. Any fixed original feasible lift remains feasible for small $\eta$
because its perturbation has order $\eta^2$ and band width order $\eta$.
The objective is uniformly coercive for small $\eta$, so minimizers
exist and comparison with this lift bounds them uniformly. A selected
cluster point satisfies all original equations and affine rows.

For a boxed optimum, bound the lifted $y$ coordinates by rational
absolute-coefficient bounds on the given $x$ box. Minimize $f+\eta P_0$
on this compact lifted polyhedron with the bands. An original optimizer
remains feasible eventually. Uniform convergence of objectives and
comparison with it make every selected cluster an original optimizer.
These supplied box rows are counted in the input.

In either case retain a fixed active affine chart and oriented support
along one convergent subsequence. A zero-dimensional chart is a point in
$K$ of polynomial encoding length. Otherwise the stationarity reduction
\eqref{eq:ncq-jacobian} and Lemma~\ref{lem:qc-elimination} give all scalar
and linear-form bounds~\eqref{eq:ncq-elimination-budget} along this one
tuple. Lemma~\ref{lem:ncq-common-field} gives the common field. The value
has absolute degree at most $D\Lambda$ and the just-proved height bound,
so its primitive integer minimal polynomial has the asserted size.
Cauchy's bound gives the coordinate radius. Remove any initial zero
coefficients from a value annihilator and apply Cauchy's bound to its
reciprocal polynomial to obtain the nonzero-value gap.
\end{proof}

\subsection{Finite infima and attained minimum-norm optimizers}
\label{app:ncq-infimum}

\begin{theorem}[Finite-infimum and attained-point encodings]
\label{thm:ncq-infimum}
For the rational system~\eqref{eq:ncq-system}, assume $S\ne\varnothing$
and $\theta=\inf_Sf> -\infty$. A nonzero integer polynomial annihilates
$\theta$, with degree and coefficient bits $L^{O(h+1)}$.
If $\theta$ is attained, at least one optimizer of least Euclidean norm
in the original coordinates has a common-field representation of degree
and total length $L^{O(h+1)}$. Its value belongs to that field, and its
squared norm has an annihilator with the same bounds. All bounds have the
coefficient-sensitive form~\eqref{eq:ncq-budget}. No supplied box is
required.
\end{theorem}

\begin{proof}
For fixed $\varepsilon>0$, put
$f_\varepsilon=f+\varepsilon\norm x^2$ and
$v_\varepsilon=\min_Sf_\varepsilon$. Since
$f_\varepsilon\ge\theta+\varepsilon\norm x^2$ on the nonempty closed
set $S$, a nonempty sublevel below its value at a feasible anchor is
compact. Hence $v_\varepsilon$ is attained. Every exact minimizer obeys,
for one fixed feasible $\bar x$,
\begin{equation}
\label{eq:ncq-anchor-bound}
 \norm x^2\le\norm{\bar x}^2+
                     \frac{f(\bar x)-\theta}{\varepsilon}.
\end{equation}
It follows that all exact minimizers and their lifts are bounded for each
fixed $\varepsilon$. Choose an unknown finite integer radius
$R_\varepsilon$ strictly enclosing all these lifts. Its magnitude may
grow as $\varepsilon\downarrow0$ and is not assigned any encoding bound.
Moreover,
\[
 \theta\le v_\varepsilon\le f(x)+\varepsilon\norm x^2
                 \quad(x\in S).
\]
Let $\varepsilon\downarrow0$ for each fixed $x$, then take the infimum
over $x$, to obtain $v_\varepsilon\to\theta$.

Choose the one perturbation tuple of Lemma~\ref{lem:ncq-perturbation}
using only original charts. Fix an admissible $\varepsilon>0$ and
minimize~\eqref{eq:ncq-perturbation} inside
$P\cap[-R_\varepsilon,R_\varepsilon]^{n+h}$. An exact regularized
optimizer remains in the bands for sufficiently small $\eta$, so these
compact problems are nonempty. Every cluster of their minimizers as
$\eta\downarrow0$ satisfies $F_j=0$. Uniform objective convergence on
this one box, and comparison with an exact regularized optimizer, show
that the cluster has $f_\varepsilon$-value $v_\varepsilon$.
Every such cluster is therefore an original global regularized minimizer
and is strictly inside the box.

Consequently all perturbed minimizers are box-interior for sufficiently
small $\eta$ at this fixed $\varepsilon$. Otherwise a sequence of
boundary minimizers would have a boundary cluster, a contradiction. Thus
the box rows disappear before the chart and KKT coefficients are formed.
No $\log R_\varepsilon$ enters their height bounds.

For each admissible $\varepsilon_\nu\downarrow0$, retain an inner
$\eta_{\nu,\mu}\downarrow0$ sequence with one chart and oriented
support. A second finite pigeonhole selection makes that label constant
on an infinite outer subsequence. The family of labels is independent of
all radii. Use the perturbed objective itself as the rational output.
Its inner limit is $v_{\varepsilon_\nu}$ and its finite outer limit is
$\theta$. Equations~\eqref{eq:ncq-jacobian} and
\eqref{eq:ncq-elimination-budget}, over $\Q$, give the required
annihilator by Lemma~\ref{lem:qc-elimination}. This takes the inner
$\eta$ limit first and the outer $\varepsilon$ limit second. Primal
points and multipliers may diverge in the outer limit. A zero-dimensional
selected chart is one fixed rational point $a$; its inner value is
$f(a_x)+\varepsilon\norm{a_x}^2$ and its outer value is rational,
so that case needs no elimination.

Now suppose attainment holds. The closed optimal set has an optimizer
$\bar x$ of minimum norm $c$: restrict a norm-minimizing sequence to
the closed ball of any optimizer. Comparison with $\bar x$ gives, for
every exact regularized minimizer,
\begin{equation}
\label{eq:ncq-minnorm-bound}
 \norm{x_\varepsilon}\le c,\qquad
 0\le f(x_\varepsilon)-\theta\le\varepsilon c^2.
\end{equation}
All their lifts now lie in one compact set. Choose one unknown box
strictly enclosing it and apply the same eventual-inactivity argument.
Select inner convergent perturbed tuples and then an outer convergent
subsequence, retaining one chart and support:
\[
 w_{\nu,\mu}\longrightarrow w_\nu\quad(\mu\to\infty),
 \qquad w_\nu\longrightarrow w^*\quad(\nu\to\infty).
\]
Inequality~\eqref{eq:ncq-minnorm-bound} applies to the inner limits and
shows that $x^*$ is an optimizer of norm $c$. Every coordinate, every
fixed rational linear form, and $\norm{x^*}^2$ uses this same nested
tuple sequence. Elimination and Lemma~\ref{lem:ncq-common-field} give
the common-field and squared-norm bounds. Since $f$ is rational,
$\theta=f(x^*)$ lies in the same field.
\end{proof}

The minimum-norm optimizer need not be unique: minimizing the zero
objective on $x^2=1$ gives two such points. We call it canonical only when
uniqueness is established. Finite infimum also does not imply attainment:
$xy\ge1$, $x\ge0$, with objective $x$, has infimum zero and no optimizer.
The theorem asserts a finite scalar limit in this case, not a limiting
feasible primal point. On a supplied nonempty rational box, compactness
gives attainment and the theorem supplies a value and an optimal point
with the same encoding bounds, counting the box in $L$.

\begin{theorem}[Nonconvex status and output with an NP oracle]
\label{thm:ncq-oracle}
For fixed $h$, an $\mathrm{FP}^{\NP}$ algorithm classifies a rational
quadratic problem as infeasible, unbounded below, finite and unattained,
or finite and attained. It returns the exact finite value, and in the
attained case can return a minimum-norm optimizer in common-field format.
The same conclusion holds for total-degree-two mixed-integer input with
explicit finite integer bounds, using the span of the continuous Hessian
blocks. The integer dimension need not be fixed in this bounded model.
\end{theorem}

\begin{proof}
Feasibility with rational threshold $f(x)\le t$ is in $\NP$ by
Theorem~\ref{thm:ncq-feasible}, with span at most $h+1$.
First test original feasibility. An effective coefficient bound from
Theorem~\ref{thm:ncq-infimum} gives $M=2^{H+2}$ such that every finite
infimum of the instance has modulus less than $M$. On a nonempty set,
the threshold $f\le-M-1$ is feasible exactly when the problem is
unbounded below. For a finite infimum, bisect the initial interval
$[-M-1,M+1]$ using threshold queries. Each answer preserves containment
of $\theta$, including a negative answer when the midpoint equals an
unattained infimum. After polynomially many queries, the degree and
coefficient bounds give the certified precision required by
Theorem~\ref{thm:qc-kll}. Recognize the primitive minimal polynomial and
refine further below its root-separation bound to isolate the correct
real root. Query lengths and all deterministic work are polynomial for
fixed $h$.

Let $P_\theta$ and $(a,b)$ encode this actual value, with nonroot
endpoints. Given a candidate common-field point, compute $\gamma=f(x)$
in its own field and verify
\begin{equation}
\label{eq:ncq-value-verifier}
 P_\theta(\gamma)=0,\qquad a<\gamma<b,
\end{equation}
as well as all original row signs. These conditions are exactly
$x\in S$ and $f(x)=\theta$; no compositum with the separately encoded
value is needed. Ask whether this verifier accepts a certificate within
the optimizer length bound of Theorem~\ref{thm:ncq-infimum}. The answer
is yes exactly when attainment holds. Standard bit-prefix search on a
padded self-delimiting bounded-length encoding returns such an optimizer
using polynomially many $\NP$ queries.

For minimum norm, Cauchy's bound gives an explicit
$B_0=2^{L^{O(h+1)}}$ containing at least one minimum-norm optimizer.
Let $\rho$ be its squared norm. Its annihilator has the theorem's bounds.
The $\NP$ question asking for a certificate of the universal minimum-norm
length bound satisfying~\eqref{eq:ncq-value-verifier} and
$\norm x^2\le t$ answers yes exactly when $t\ge\rho$. Bisect
$[0,nB_0^2]$ and recognize $\rho$ by Theorem~\ref{thm:qc-kll}; this
requires interval containment, not an infeasible lower endpoint when
$\rho=0$. Prefix search with both exact scalar equalities returns one
minimum-norm optimizer. Equality to $\rho$ is tested in the candidate's
field using its polynomial and isolator, as for $\theta$.

For bounded integer variables, guess their polynomial-length assignment
and use the continuous witness for that slice. Substitution preserves the
continuous Hessian span and gives uniform polynomial coefficient bounds.
There are finitely many assignments. A feasible unbounded slice makes the
mixed problem unbounded; otherwise its finite infimum equals one of the
finite slice infima. Attainment has a witness in a slice attaining the
finite mixed value. The least squared norm in the original mixed
coordinates is the minimum, over the finitely many such slices, of
$\norm z^2$ plus the least continuous squared norm. One slice attains
this minimum, and its scalar bound includes the constant $\norm z^2$.
The uniform point and scalar bounds therefore
justify the same queries, recognition, and prefix search.
\end{proof}

The represented point verifies feasibility. Its global optimality in this
algorithm follows from the separately computed exact value; no short
standalone certificate of nonconvex global optimality is asserted.
Neither the output bounds nor the oracle theorem imply an ordinary
polynomial-time global optimization algorithm. Indeed, on $[0,1]^n$ the
single concave row $\sum_i x_i(1-x_i)\le0$ forces every coordinate to
be Boolean. Affine three-literal clause inequalities then give strong
$\NP$-hardness of feasibility with $h=1$. Irrational feasible points
are already necessary for the span-one pair $x^2\le2$, $2-x^2\le0$.

Attainment remains strongly $\NP$-hard with one nonlinear row, a linear
objective, and promised known finite infimum zero. For a three-literal
formula, impose its affine clause rows on $x\in[0,1]^n$, add $t,y\ge0$,
and minimize $t$ subject to
\[
 \sum_i x_i(1-x_i)-ty\le0.
\]
Every $t>0$ is feasible with $x_i=1/2$ and $y=n/(4t)$, so the infimum
is zero. At $t=0$ the row forces a Boolean satisfying assignment, and any
such assignment with $y=0$ attains zero. The unique native Hessian
direction is nonzero and all coefficients are bounded. On this promised
subclass attainment is in $\NP$ by appending $t\le0$.
These bounded-integer statements make no claim about arbitrary unbounded
integer domains; Proposition~\ref{prop:qc-pell} explains an expanded-output
obstruction in that setting.

\subsection{Cone input over one selected real field}
\label{app:cone-setting}

Represent
\begin{equation}
\label{eq:cone-field}
 K=\Q(\alpha)\subset\R,\qquad p(\alpha)=0,\qquad a<\alpha<b,
\end{equation}
by a dense primitive irreducible integer polynomial $p$ and rational
nonroot endpoints isolating its selected real root. All coefficients are
explicit rational power-basis vectors of length $D=\deg p$. Their
vectors, all matrix entries, the dense $p$, and the isolator are counted
in $L$, so $D\le L$. The degree varies with the input. A succinct
high-degree field encoding or separately encoded unrelated algebraic
coefficients is outside this model. A real field here means a chosen real
embedding, not a field with only one real embedding.

For $w=(z,x)\in\Z^t\times\R^n$, impose affine weak rows and equations
over $K$ and
\begin{equation}
\label{eq:cone-system}
 \norm{u_i(w)}\le t_i(w),\qquad
 u_i(w)=A_iw+b_i,\quad t_i(w)=c_i^{\mathsf T}w+d_i.
\end{equation}
The scalar quadratic residual is
$q_i=\norm{u_i}^2-t_i^2$. Its continuous Hessian is
\begin{equation}
\label{eq:cone-span}
 H_i=2(A_{ix}^{\mathsf T}A_{ix}-c_{ix}c_{ix}^{\mathsf T}),
 \qquad h=\dim_K\Span_K\{H_i:i\le m\}.
\end{equation}
The squared description always retains $t_i\ge0$. Flattening the
matrices reduces their span to a rank over $K$; nonzero minors show that
this equals their real span rank at the selected embedding. It need not
equal their rational span: $H$ and $\sqrt2H$ have $K$-span one and
rational span two when $H\ne0$ is rational. These matrices can be
indefinite even though the original conic set is convex.

\begin{theorem}[Exact common-field cone feasibility]
\label{thm:cone-feasibility}
For each fixed $t,h$, feasibility of~\eqref{eq:cone-system} with
arbitrary affine rows over~\eqref{eq:cone-field} is decidable in
polynomial Turing time in $L$. If feasible, a polynomial-length integer
assignment $z$ whose original exact continuous fiber is nonempty can be
returned. No boxes, Slater condition, full dimensionality, or closedness
of the projection onto $z$ are assumed. The polynomial exponent may
depend on $t,h$.
\end{theorem}

We prove the theorem by first establishing a small integer assignment and
a uniform exact residual gap, and then constructing a rational polyhedron
that preserves precisely the feasible integer assignments within the
resulting boxes. The integer-witness theorem is used only for a radius;
the large formula in that proof is never constructed by the algorithm.

\subsection{An exact compressed projection and its integer radius}
\label{app:cone-projection}

Let $Y\subset\R^t$ be the real projection of the original conic set.
It is convex and may be nonclosed. First apply the lift
\eqref{eq:ncq-lift} only to the continuous Hessians of the $q_i$. Writing
$v=(x,y)$ gives
\begin{equation}
\label{eq:cone-fiber-lift}
 v\in P_z,\qquad F_j(v)=\tfrac12x^{\mathsf T}B_jx-y_j=0.
\end{equation}
The row normals of $P_z$ are polynomials in $z$ of degree at most one,
and their constant terms have degree at most two. Every affine row and
cone sign is retained.

For each possible active affine subset, include every pivot-row and
pivot-column choice in its equality system $C(z)v=d(z)$. On a nonzero
pivot minor $\sigma(z)$, Cramer's rule gives
$v=a(z)+V(z)u$. Include guards
\begin{equation}
\label{eq:cone-chart-guards}
 \sigma(z)\ne0,\qquad C(z)a(z)=d(z),\qquad C(z)V(z)=0.
\end{equation}
The free rows of $V$ form the identity, so these guards assert that the
chart parametrizes the whole equality solution space, including nonpivot
equations. The rank-zero chart uses $\sigma=1$, $a=0$, $V=I$.
All chart numerators and denominators have polynomial degrees and
coefficient bit lengths over $K$. There are at most exponentially many
charts. Include also their zero-dimensional cases.

Fix any real $z$. On every valid positive-dimensional chart use the
one-parameter objective $\norm v^2+\varepsilon P_0(v)$ and bands
$|F_j+\varepsilon^2P_j|\le\varepsilon$. The genericity argument and
integer grid of Lemma~\ref{lem:ncq-perturbation} give a good choice on
all valid charts of this fixed $z$. The degree and chart-count bounds
are uniform in $z$, regardless of the specialized coefficients' heights.
Thus one common finite grid $\mathcal G$ of polynomial-bit integer
perturbations contains a good choice for every $z$. The selected grid
point may depend on $z$; no single perturbation generic for every $z$
is claimed. Include every grid point as a Boolean disjunct.

On a chart of dimension $d>0$ and a selected oriented support of size
$s\le\min(h,d)$, write stationarity with $\lambda\ge0$ and eliminate
$u$ by the adjugate formula on the guard $\det M\ne0$. This gives a
rational-function point $v_*(z,\varepsilon,\lambda)$. Require all
guards~\eqref{eq:cone-chart-guards}, its common denominator nonzero,
$\lambda\ge0$, $0<\varepsilon<\delta$, and
\begin{equation}
\label{eq:cone-projection-atoms}
 \norm{v_*}^2\le R,\quad v_*\in P_z,\quad
 |F_j(v_*)+\varepsilon^2P_j(v_*)|\le\varepsilon,
\end{equation}
together with equality on each selected oriented band. Clear denominators
by even powers under their nonzero guards; no pivot sign is presumed.
Polynomial-size determinants and a fixed number of substitutions give
polynomial degree and individual coefficient bits for every atom.
A zero-dimensional disjunct directly requires that its point belong to
$P_z\cap Z(F_1,\ldots,F_h)$ and have squared norm at most $R$.
Pad multipliers to length $h$.

The exact projection formula over $K$ has prefix
\begin{equation}
\label{eq:cone-projection-prefix}
 \exists R>0\quad\forall\delta>0\quad
       \exists(\varepsilon,\lambda_1,\ldots,\lambda_h).
\end{equation}
If $z\in Y$, fix one feasible lift. A good grid point makes it feasible
in the outward bands eventually. The objective is uniformly coercive,
so global minimizers exist and have bounded norm by comparison with that
lift. Their active charts and band supports have at most $h$ nonlinear
rows, and genericity gives the required KKT multipliers and invertible
$M$. Hence for every $\delta>0$ a sufficiently small $\varepsilon$
and such a minimizer satisfy a disjunct with one fixed finite $R$.
If zero-dimensional charts occur on a sequence tending to zero, one
constant chart gives a directly feasible limiting point instead.

Conversely, suppose the formula holds at a fixed $z$. A direct
zero-dimensional disjunct already supplies a feasible point. Otherwise
take $\delta=1/j$. The corresponding points have squared norm at most
the one fixed $R$ and $\varepsilon\to0$. They have a convergent
subsequence. The finite grid has bounded coefficients, so its perturbing
polynomials are bounded on this ball even if their choice changes.
Equation~\eqref{eq:cone-projection-atoms} therefore gives $F_j\to0$;
all original affine rows persist by closedness of the fixed $P_z$. The
limit satisfies~\eqref{eq:cone-fiber-lift}. This proves exact projection,
including at nonclosed $Y$. Moving $R$ inside the universal block would
allow escaping witnesses and would invalidate the reverse implication.
The affine case $h=0$ is captured by the minimum-norm active affine
chart with its invertible restricted norm Hessian.

Now reduce all field arithmetic modulo $p$ and replace each coefficient
by its rational polynomial in one new real variable $v_0$. Conjoin
$p(v_0)=0$, $a<v_0<b$, and put $v_0$ with $R$ in the first block.
Clear remaining rational coefficient denominators positively. The
isolator selects exactly the intended embedding. Atomic integer degrees
and individual coefficient bits stay polynomial in $L$, since $D\le L$.
The three quantified block dimensions are
\begin{equation}
\label{eq:cone-blocks}
 2,\qquad1,\qquad h+1.
\end{equation}
There is one generator coordinate, not $D$ quantified coordinates.

For the radius proof, first adjoin one free coordinate fixed to zero,
giving the convex set $\{0\}\times Y$ in $t+1$ free variables. Apply
fixed-block quantifier elimination
\cite[Theorem~2.27]{Basu2014} to this mathematical description. For
input degree bound $d_*\ge2$, integer coefficient-bit bound $b_*$, quantified
block sizes $n_1,\ldots,n_\omega$, and $\ell$ free variables, that
theorem bounds each output polynomial's degree by
$d_*^{O(n_\omega)\cdots O(n_1)}$ and its coefficient bits by
$b_*d_*^{O(n_\omega)\cdots O(n_1)O(\ell)}$.
These per-polynomial bounds are independent of the number of input
atoms; output count and running time can depend on it. The
blocks~\eqref{eq:cone-blocks} and $\ell=t+1$ therefore give an
equivalent quantifier-free formula with individual degrees and bits
\[
 d'=L^{O(h+1)},\qquad b'=L^{O((h+1)(t+1))}.
\]
Neither this formula nor the elimination is constructed by our algorithm.

Now use only the quantifier-free case of Khachiyan--Porkolab's
integer-witness theorem \cite[Theorem~1.1]{KhachiyanPorkolab2000}.
It allows arbitrary Boolean formulas and nonclosed or lower-dimensional
convex sets. If an integer optimum exists in free dimension $t+1$, an
optimal integer point has coordinate bit bound
$b'(d')^{O((t+1)^4)}$, independently of the number of atoms. Minimizing
the adjoined coordinate fixed to zero is precisely their reduction from
feasibility to optimization. The preceding degree and bit bounds give an
effective $B_z=L^{C(h+1)(t+1)^4}$ such that
\begin{equation}
\label{eq:cone-integer-radius}
 Y\cap\Z^t\ne\varnothing\ \Longrightarrow\
 \exists z\in Y\cap\Z^t:\quad\norm z_\infty\le2^{B_z}.
\end{equation}
Only this atom-count-independent witness bound is used. The theorem's
formula-size-dependent algorithm is not applied to either large formula.
When $t=0$ this integer-radius step is unnecessary.

Print the integer box $[-Z,Z]^t$, $Z=2^{B_z}$. Every integer substitution
in it has uniformly polynomial encoding length for fixed $t,h$.
Lemma~\ref{lem:ncq-field-bounds} then gives one rational continuous box
$[-R_x,R_x]^n$ meeting every nonempty exact fiber in that integer box.
Its endpoint bit lengths are polynomial for fixed $t,h$. No fiber or
integer assignment is enumerated. These boxes preserve existence; they
need not contain all feasible points or assignments.

\subsection{A uniform gap and a rational linear reduction}
\label{app:cone-gap}

For boxed integer $z$ define
\begin{equation}
\label{eq:cone-residual-value}
 \gamma_z=\min_{x\in[-R_x,R_x]^n}
 \max\{0,q_i(z,x),-t_i(z,x),\ell_j(z,x),e_j(z,x),-e_j(z,x)\},
\end{equation}
where $\ell_j\le0$ and $e_j=0$ list all original affine rows.
The maximum ranges over every indicated scalar row. It is continuous
on a compact box, so the minimum exists and is zero exactly when the
original boxed fiber is feasible. Its epigraph uses one new variable,
which adds no Hessian direction. A direct absolute-coefficient bound on
the box gives a rational upper bound making that epigraph compact and
nonempty. Lemma~\ref{lem:ncq-field-bounds} supplies one effective
\begin{equation}
\label{eq:cone-gap}
 \Delta=2^{-B}\le1,\qquad B=L^{C_{t,h}},\qquad
                \gamma_z=0\ \text{or}\ \gamma_z\ge\Delta
\end{equation}
for every boxed integer assignment. This gap belongs to the original
exact system.

We record an explicit rational cone approximation to make its bit
contract and apex behavior precise. The construction is the
Pythagorean-triple approximation and binary tree of Kocuk
\cite[Sections~2.3 and 3.3]{Kocuk2021}, following the lifted cone
approximation of Ben-Tal--Nemirovski~\cite{BenTalNemirovski2001}.

\begin{lemma}[Rational Lorentz lift]
\label{lem:cone-rational-lift}
For rational $0<\epsilon\le1$, a homogeneous rational linear lift of
polynomial size and coefficient bits in $d+\log(1/\epsilon)$ has
projection $P_{d,\epsilon}$ satisfying
\[
 \{(u,s):\norm u\le s\}\subseteq P_{d,\epsilon}
 \subseteq\{(u,s):s\ge0,\ \norm u\le(1+\epsilon)s\}.
\]
It is constructible in polynomial bit time, including at $s=0$.
\end{lemma}

\begin{proof}
Dimensions $d=0,1$ are exactly linear. For two coordinates use the
integer triples
\begin{align*}
 (a_1,b_1,c_1)&=(120,119,169),\qquad k_j=2^{j-2}+2,\\
 (a_j,b_j,c_j)&=(2k_j-1,2k_j^2-2k_j,2k_j^2-2k_j+1)
                  \quad(j\ge2).
\end{align*}
For tolerance $\delta>0$ choose the first $J\ge2$ with
$\delta b_J\ge1$. Introduce nonnegative folding coordinates by
\begin{align*}
 &\xi_0\ge\pm u_1,\qquad\eta_0\ge\pm u_2,\\
 &c_j\xi_j=b_j\xi_{j-1}+a_j\eta_{j-1},\qquad
 c_j\eta_j\ge\pm(-a_j\xi_{j-1}+b_j\eta_{j-1})\quad(1\le j\le J),\\
 &\xi_J\le s,\qquad b_J\eta_J\le a_J\xi_J.
\end{align*}
Each $\pm$ means both inequalities. The triples satisfy
$a_j^2+b_j^2=c_j^2$. If $\theta_j$ is the first-quadrant angle with
sine $a_j/c_j$, then $\theta_1\ge\pi/4$,
$\theta_2=\theta_1/2$, and
$\theta_{j-1}/2\le\theta_j\le\theta_{j-1}$ for $j\ge3$.
For the lower comparison write $k=k_{j-1}$ and $k_j=2k-2$;
the half-angle comparison reduces to
$(4k-5)/(8k^2-20k+12)\ge1/(2k-1)$, whose cross-product difference
is $6k-7>0$. The upper comparison follows since
$(2k-1)/(2k^2-2k)$ decreases for $k\ge2$.
Successive exact rotations and reflections therefore fold the first
quadrant into angles $[0,\theta_j]$. Taking every absolute-value slack
at equality lifts each exact two-coordinate cone point.

Conversely the folding coordinates are nonnegative; rotations preserve
norm before slack is added, so
\[
 \sqrt{u_1^2+u_2^2}\le\sqrt{\xi_J^2+\eta_J^2}
       \le(c_J/b_J)\xi_J\le(1+\delta)s.
\]
This also proves $s\ge0$ and excludes a spurious branch at the apex.
For $j\ge2$, $c_j=b_j+1$ and $b_{j+1}<4b_j$. Minimality of $J$
bounds every coefficient by $\max(169,4/\delta+1)$, including $J=2$.
Since $b_j\ge2^{2j-3}$, the exact rational stopping test uses
$O(1+\log^+(1/\delta))$ stages.

For general $d\ge2$, pad with zeros to $2^k$ leaves, where
$k=\lceil\log_2d\rceil$, and use this lift on each internal vertex
of a balanced binary tree with $\delta=\epsilon/(2k)$.
There are fewer than $2d$ internal vertices. True subtree norms give inner
inclusion, and induction gives outer factor $(1+\delta)^k$.
The inequalities $\log(1+\epsilon)\ge\epsilon/2$ and
$(1+\epsilon/(2k))^k\le\exp(\epsilon/2)$ give factor at most
$1+\epsilon$. Coefficient generation and stopping use integer and
rational arithmetic only, with the asserted size and bit bounds.
\end{proof}

Choose a rational $T\ge1$ bounding both $\norm{u_i(w)}$ and
$|t_i(w)|$ on the printed boxes. Absolute-coefficient sums give such
a $T$ with polynomial bits. Set
\begin{equation}
\label{eq:cone-tolerances}
 \eta=\frac{\Delta}{256T},\qquad
 \epsilon=\frac{\Delta}{256T^2}.
\end{equation}
Approximate every affine map rationally so that, uniformly on the boxes,
\[
 \norm{u_i-\widehat u_i}\le\eta,\quad
 |t_i-\widehat t_i|\le\eta,\quad
 |\ell_j-\widehat\ell_j|\le\eta,\quad
 |e_j-\widehat e_j|\le\eta.
\]
For example, with $V_0=\max(Z,R_x,1)$ and $d_0$ bounding all cone
vector dimensions and one, coefficient error at most
$\eta/((t+n+1)V_0(d_0+1))$ suffices, using the vector $\ell_1$ norm.
This accuracy is computable in polynomial time. On the input isolator a
coefficient polynomial $g(v)=\sum_ra_rv^r$ has derivative bounded by
$\sum_{r\ge1}r|a_r|H_0^{r-1}$, with
$H_0=\max(1,|a|,|b|)$. This bound has polynomial bit length, so dense
root refinement for $\alpha$ to a polynomial number of bits controls
all evaluations even if the derivative magnitude is large.

Use rational affine rows $\widehat\ell_j\le\eta$ and
$|\widehat e_j|\le\eta$, and apply Lemma~\ref{lem:cone-rational-lift}
to $(\widehat u_i,s_i)$ with $s_i=\widehat t_i+2\eta$. Keep both
boxes exact. Every original feasible point lifts, since
$\norm{\widehat u_i}\le\norm{u_i}+\eta\le s_i$.
Conversely a lifted point has $s_i\ge0$,
$\norm{\widehat u_i}\le(1+\epsilon)s_i$, original affine residuals
at most $2\eta$, and $-t_i\le3\eta$. Its squared residual obeys
\begin{align}
\label{eq:cone-rounding-error}
 q_i&\le3\epsilon(T+3\eta)^2+
                 \eta(2T+\eta)+3\eta(2T+3\eta)\notag\\
    &\le48\epsilon T^2+18T\eta
       =\frac{66}{256}\Delta<\Delta/2.
\end{align}
The first line uses the absolute error of squared norms and squared right
sides, without presuming the true $t_i$ nonnegative. For the second line
use $\epsilon,\eta\le1$ and $T\ge1$. The sign residual bound is
therefore essential. All residuals in~\eqref{eq:cone-residual-value}
are less than $\Delta/2$. At integral $z$, the exact
gap~\eqref{eq:cone-gap} forces $\gamma_z=0$.

The resulting linear system is rational, of polynomial length for fixed
$t,h$, and has only the original $t$ integer variables. Rational mixed
integer linear feasibility with fixed integer dimension and arbitrary
continuous dimension is polynomial-time and returns an integer assignment
\cite[Section~5]{Lenstra1983}. Applying it proves
Theorem~\ref{thm:cone-feasibility}. The linear lift's displayed continuous
point need not satisfy the original cones. Rational rounding can increase
the squared-Hessian span; no span-dependent theorem is applied to the
rounded system. The radius and residual gap were established for the
original exact system.

\subsection{The absolute field of a canonical continuous witness}
\label{app:cone-witness}

\begin{theorem}[Exact common-field cone witnesses]
\label{thm:cone-witness}
For a nonempty continuous system~\eqref{eq:cone-system} over
\eqref{eq:cone-field}, its unique minimum-norm feasible point $x^*$
has absolute field $E=K(x^*)$ of degree and total common-field
representation length $L^{O(h+1)}$. For fixed $h$, an algorithm running
in $L^{C_h}$ bit operations decides feasibility and, if feasible, returns
an irreducible primitive polynomial $P$, a selected real root $\beta$,
and rational coordinate maps
\[
 \alpha=b_0(\beta),\qquad x_j^*=b_j(\beta),\qquad\Q(\beta)=E.
\]
For fixed $t,h$, a mixed-integer system likewise admits polynomial-time
recovery of a full feasible point: first return $z$ and then the unique
minimum-norm continuous point in its original exact fiber.
\end{theorem}

\begin{proof}
The continuous conic set $C$ is closed and convex. Coercivity of the
squared norm gives existence and strict convexity gives uniqueness of
$x^*$. Lift it to $w^*$ and choose an unknown integer box strictly
enclosing that lift. On the compact intersection of this box with the
original lifted polyhedron, minimize $\norm x^2+\eta P_0(w)$ subject
to the generic outward bands. The point $w^*$ remains feasible eventually.
Every cluster of minimizers is exactly feasible and, by uniform objective
convergence and comparison with $w^*$, has squared norm at most
$\norm{x^*}^2$. Uniqueness of $x^*$ and its lift forces every cluster
to be $w^*$. Thus every auxiliary box row is inactive eventually, by the
same boundary-sequence contradiction used in
Theorem~\ref{thm:ncq-infimum}.

Retain one original affine chart and oriented support along a sequence.
The nonconvex KKT reduction has $s\le h$ nonsingular multiplier
variables and no box coefficients. Every fixed $K$-linear form in $x^*$
is an output on this same sequence, of relative degree at most
$\Lambda=L^{O(h+1)}$. Lemma~\ref{lem:ncq-common-field} proves
$[K(x^*):K]\le\Lambda$, hence $[E:\Q]\le D\Lambda$, and the
absolute coefficient and representation bounds. In particular an effective
Cauchy bound prints a rational box $[-B,B]^n$ containing $x^*$ itself.
A box merely meeting $C$ would not justify this step.

Decide feasibility by Theorem~\ref{thm:cone-feasibility} with $t=0$.
For an approximation request $q\ge1$, set $\tau=2^{-q}$ and restart
from $[-B,B]^n$. For rational $r\ge0$ use
\begin{equation}
\label{eq:cone-norm-query}
 \norm x^2\le r\quad\Longleftrightarrow\quad
                   \norm{(2x,r-1)}\le r+1.
\end{equation}
The squared residual is $4\norm x^2-4r$ and its Hessian is $8I$, so
the feasibility queries have span at most $h+1$ and use only the
original field. Bisect $[0,nB^2]$ until a feasible rational upper
threshold $u$ satisfies
$\norm{x^*}^2\le u\le\norm{x^*}^2+\tau^2/16$.
The projection inequality for a closed convex set gives, for every
$x\in C$,
\[
 \ip{x^*}{x-x^*}\ge0,\qquad
 \norm{x-x^*}^2\le\norm x^2-\norm{x^*}^2.
\]
Hence all points of the nonempty norm slice are within $\tau/4$ of
$x^*$. Bisect closed coordinate intervals while preserving nonempty
intersection with that slice: retain the lower half if feasible, and
otherwise the upper half. After all widths are at most $\tau$, the
rational box midpoint is within $3\tau/4$ of each coordinate of
$x^*$. It need not itself be feasible. Restarting the box for each
accuracy request ensures approximation of this same specified point.

Refine the input isolator independently to approximate $\alpha$. This
gives a certified oracle for the tuple $(\alpha,x^*)$, with polynomial
query number and length in the input, printed radius bits, and $q$.
Apply Theorem~\ref{thm:qc-recovery} with joint degree $D\Lambda$ and
the absolute coordinate coefficient bounds. Its recognition precision is
polynomial in these quantities; it constructs the common field with
polynomial overhead. Including $\alpha$ in the tuple gives exactly the
map $b_0$ and preserves the selected embedding. The printed radius and
precision enter later feasibility calls, so composing their polynomial
bounds gives $L^{C_h}$ time; we claim no sharper exponent for recovery.

Verify $p(b_0(\beta))=0$ and $a<b_0(\beta)<b$. Map all input
coefficients through $b_0$, reduce modulo $P$, and check every original
affine row, squared residual, and cone sign by Lemma~\ref{lem:qc-sign}.
This verifies original feasibility and the input embedding. The
minimum-norm property follows from the algorithm and its specified tuple,
not from these feasibility checks alone. If $n=0$, field sign tests decide
the scalar conditions and the output retains the input field.

For mixed integers substitute the polynomial-length $z$ returned by
Theorem~\ref{thm:cone-feasibility} into the original exact system and
apply the continuous algorithm. The field and continuous Hessians remain
unchanged and the fiber input length is polynomial for fixed $t,h$.
Its canonical point is canonical within this fiber; no global minimum-norm
integer assignment is asserted.
\end{proof}

The extension field can be proper. With $K=\Q(\sqrt2)$ in its positive
embedding, impose $y=\sqrt2$ and
\[
 \norm{(\sqrt2,1)}\le x,\qquad \norm{(x,\sqrt2x)}\le3.
\]
The first cone forces $x\ge\sqrt3$ and the second gives
$3x^2\le9$, so the unique point has $x=\sqrt3$. The squared-Hessian
span is one, while $K(x,y)=\Q(\sqrt2,\sqrt3)$ has absolute degree four.
Recovering just $\Q(x)$ would fail to retain the input field.

\subsection{Known algebraic levels and the scope of the algorithms}
\label{app:cone-levels}

\begin{corollary}[Known thresholds and attainment]
\label{cor:cone-threshold}
For rational original cone data, a rational affine $f$, and one explicitly
isolated algebraic threshold $\theta$, feasibility after appending
$f\le\theta$ and full feasible-point recovery are polynomial-time for
fixed $t,h$, including the threshold encoding in the input. The parameter
$h$ is the original rational squared-Hessian span. If $\theta$ is the
known finite infimum, this decides attainment and returns an optimizer
when attainment holds.
\end{corollary}

\begin{proof}
Work in the single selected field $\Q(\theta)$. Extending scalars
preserves the rank of the rational matrix family, and the affine level
adds no Hessian direction. Apply Theorems~\ref{thm:cone-feasibility}
and~\ref{thm:cone-witness}. If $\theta$ is the true finite infimum,
a feasible point with $f\le\theta$ exists exactly when it is attained.
The boxes and gap in the proof include the added level from the outset.
\end{proof}

\begin{corollary}[Recovery at a known fractional optimum]
\label{cor:cone-fractional}
Suppose the original cone data and affine functions $f,d$ are rational,
and an exact finite attained optimum $\theta$ of $f/d$ on the domain
$d>0$ is supplied. For fixed $t,h$, an exact optimizer is recoverable
in polynomial time in the original data plus the encoding of $\theta$.
The construction adds at most one continuous Hessian direction.
\end{corollary}

\begin{proof}
Over $\Q(\theta)$ impose $f-\theta d=0$, introduce one continuous
variable $s$, and impose the reciprocal cone
\[
 \norm{(2,d-s)}\le d+s.
\]
Its squared residual is $4-4ds$, and its sign row is $d+s\ge0$.
Thus it implies $ds\ge1$ and $d+s\ge0$, which force $d,s>0$.
Conversely $d>0$ admits $s=1/d$. The augmented system is closed and
adds at most one squared-Hessian direction. The known attainment makes
it nonempty; recover a point by the preceding theorems and discard $s$.
Its positive denominator and optimal-level equation make it an optimizer.
\end{proof}

These corollaries use a supplied value; they do not compute an unknown
mixed-integer infimum or give a general bound on its encoding. The rational
original-data restriction in the fractional corollary ensures that all
coefficients lie in the one supplied field $\Q(\theta)$.

For fixed $h$ the continuous cone decision and witness algorithms are
ordinary polynomial-time algorithms, and for fixed $t,h$ the same holds
without integer boxes. With arbitrary integer dimension and supplied
finite polynomial-bit integer bounds, the same radius, gap, and rational
lift give polynomial-size certificates and a polynomial-size rational
mixed integer linear reduction for fixed $h$; the conclusion is $\NP$
membership, not a deterministic polynomial-time algorithm. The
unbounded-integer radius~\eqref{eq:cone-integer-radius} is not polynomial
uniformly in $t$, so it gives no such $\NP$ assertion for arbitrary
unbounded integer dimension. Nor do any of these bounds imply an expanded
output runtime of the form $F(h)L^C$ with an absolute exponent: the output
lower bounds in Appendix~\ref{app:qc} already exclude that conclusion
for its covered formats.

Convexity enters the cone algorithm through the convex real projection and
the outward polyhedral cone lift. The quadratic arithmetic lemmas alone
are not efficient nonconvex optimization methods. The input field is one
explicit selected field, and its absolute output extension is constructed
only after approximation of one fixed tuple. These distinctions separate
exact decision, algebraic output length, and verification of the resulting
point.
